% ======================================================================
% Kapitel 3 — Solution: Design and Implementation (Budget: ~9 Seiten)
% Einzige Heimat fuer: Designkriterien (3.2) und Evaluationsdesign (3.8).
% Architektur begruendet nichts neu, sondern referenziert Kriterien-IDs.
% ======================================================================
\chapter{Solution: Design and Implementation}
\label{ch:solution}

\section{Task formulation}
\label{sec:task-formulation}

The task addressed in this thesis is grounded question answering over CTI
sources. The prototype operates on a corpus $\mathcal{D}$ of documents
collected from heterogeneous CTI sources and segmented into a set of
passages $\mathcal{P}$ (Section~\ref{sec:ingestion}). Given an analyst
query $q$, a retrieval function $R$ selects the $k$ highest-ranked passages
as context:
\begin{equation}
    \mathcal{C}_k = R(q, \mathcal{P}), \qquad \lvert\mathcal{C}_k\rvert = k.
    \label{eq:retrieval}
\end{equation}
A generation model $g$ then produces the answer
\begin{equation}
    a = g(q, \mathcal{C}_k),
    \label{eq:generation}
\end{equation}
subject to a grounding constraint: every factual claim in $a$ must be
attributable to at least one passage in $\mathcal{C}_k$ and marked with an
inline citation to that passage. The analyst can therefore verify each
claim against its source instead of trusting the answer as a whole
(bottleneck B3, Section~\ref{sec:cti}).

In the hybrid configuration, $R$ combines a lexical relevance score
(BM25) with a semantic similarity score (dense embeddings) and reranks the
fused candidates; Section~\ref{sec:retrieval} describes this instantiation.
Setting $\mathcal{C}_k = \emptyset$ yields the non-retrieval baseline
\begin{equation}
    a_0 = g(q, \emptyset),
    \label{eq:baseline}
\end{equation}
in which the model answers from parametric knowledge alone. This baseline
corresponds to configuration V0 in the evaluation
(Section~\ref{sec:baselines}); the comparison between $a$ and $a_0$
quantifies the effect of retrieval augmentation and thereby
operationalizes SRQ1.

\section{Requirements and design criteria}
\label{sec:requirements}
% TODO: EINZIGE Stelle fuer Designkriterien. Aus Forschungsluecken
% (\ref{sec:research-gap}) UND Workflow-Engpaessen B1--B3 (\ref{sec:cti})
% ableiten --- im Einleitungstext der Section beide Quellen explizit nennen.
% Rueckverfolgbarkeitskette: B/G -> Criterion -> Result -> RQ-Antwort.
% Zuordnung (fuer Herleitungstext): C1 <- B3/G1, C2 <- B1/G1, C3 <- G3,
% C4 <- G2, C5 <- B1/G2, C6 <- Proposal-Vorgabe (open-source only).

% Kriterien-Entwuerfe direkt aus SotA Kap. 4.2 "Implications for this
% thesis" + Proposal (open-source only). Beim Schreiben: pro Kriterium
% 1--2 Saetze Herleitung mit Gap-Bezug (G1--G3), dann nur noch per ID
% referenzieren.
\begin{table}[htbp]
    \centering
    \small
    \begin{tabular}{lp{0.62\textwidth}ll}
        \toprule
        \textbf{ID} & \textbf{Design criterion} & \textbf{Gap} & \textbf{Ref.} \\
        \midrule
        C1 & Evidence-linked outputs: explicit, source-based citations for
             key claims to support traceability and reduce verification
             effort & G1 & SRQ1 \\
        C2 & Hybrid retrieval as a baseline: combined lexical (BM25) and
             semantic retrieval with optional reranking & G1 & SRQ1 \\
        C3 & Source integrity and reliability: curated ingestion and
             source-reliability signals to limit misinformation risk & G3 & MRQ \\
        C4 & Scenario-based evaluation: assessment on representative CTI
             analyst tasks beyond offline metrics & G2 & SRQ2 \\
        C5 & Analyst-oriented presentation: structured outputs matching
             analyst intent (task coverage, clarity, transparency) & G2 & SRQ2 \\
        C6 & Open-source, locally deployable stack (data protection,
             reproducibility, cost) & --- & MRQ \\
        \bottomrule
    \end{tabular}
    \caption{Design criteria with traceability to research gaps and research questions}
    \label{tab:design-criteria}
\end{table}

\section{System architecture}
\label{sec:architecture}
% TODO: Komponenten und Datenfluss beschreiben. Designentscheidungen NUR
% per Kriterien-ID begruenden ("addresses C2"), keine erneute Herleitung.
...

% TODO Abbildung: Export aus docs/architecture/cti_rag_architecture_overview.drawio
% -> images/ (PDF/PNG). Konkrete Instanz von \ref{fig:rag-pipeline}.
% \begin{figure}[htbp]
%     \centering
%     \includegraphics[width=\textwidth]{cti_rag_architecture_overview}
%     \caption{System architecture of the CTI-RAG prototype}
%     \label{fig:architecture}
% \end{figure}

\section{Data ingestion and preprocessing}
\label{sec:ingestion}
% TODO: Quellenkorpus (welche CTI-Quellen, Umfang), Parsing, Chunking,
% Metadaten.
...

% TODO Tabelle: Korpus-Statistik aus dem Index ziehen.
% \begin{table}[htbp]
%     \centering\small
%     \begin{tabular}{llrr}
%         \toprule
%         \textbf{Source} & \textbf{Document type} & \textbf{Documents} & \textbf{Chunks} \\
%         \midrule
%         CVE/NVD & ... & ... & ... \\
%         CISA advisories & ... & ... & ... \\
%         MISP & ... & ... & ... \\
%         Vendor reports & ... & ... & ... \\
%         \bottomrule
%     \end{tabular}
%     \caption{Corpus statistics of the ingested CTI sources}
%     \label{tab:corpus-stats}
% \end{table}

\section{Hybrid retrieval and reranking}
\label{sec:retrieval}
% TODO: dense + BM25, Fusion, Reranker. Kein eigenes Diagramm ---
% der Fluss steckt in \ref{fig:architecture}.
...

% TODO Tabelle: Werte aus configs/settings.yaml uebernehmen.
% \begin{table}[htbp]
%     \centering\small
%     \begin{tabular}{llp{0.42\textwidth}}
%         \toprule
%         \textbf{Component} & \textbf{Model / method} & \textbf{Parameters} \\
%         \midrule
%         Sparse retrieval & BM25 & ... \\
%         Dense retrieval & <embedding model> & ... \\
%         Fusion & ... & ... \\
%         Reranking & <cross-encoder> & top-k = ... \\
%         \bottomrule
%     \end{tabular}
%     \caption{Retrieval pipeline configuration}
%     \label{tab:retrieval-config}
% \end{table}

\section{Context assembly, answer generation and citation}
\label{sec:generation}
% TODO: Kontextaufbau, Prompt-Struktur (Verweis auf Anhang Prompt Templates),
% Grounding-/Zitiermechanismus, Refusal-Verhalten bei fehlender Evidenz.
...

% TODO Listing: gekuerztes System-Prompt mit Zitier-Instruktion
% (aus src/cti_rag/rag/prompts.py); Vollversion in Anhang B.
% \begin{lstlisting}[caption={System prompt with citation instruction (excerpt)},
%                    label={lst:system-prompt}]
% ...
% \end{lstlisting}

\section{Implementation}
\label{sec:implementation}
% TODO: Technologie-Stack, Verweis auf Repository (Anhang).
% Hinweis: MCP-/Agent-Erweiterung hier nur kurz erwaehnen oder ganz in
% Future Work (\ref{sec:future-work}) verschieben.
...

% TODO Tabelle: Komponenten nach Salmas Vorbild (ihre Tabelle 11).
% \begin{table}[htbp]
%     \centering\small
%     \begin{tabular}{lp{0.22\textwidth}p{0.45\textwidth}}
%         \toprule
%         \textbf{Path} & \textbf{Function} & \textbf{Details} \\
%         \midrule
%         src/cti\_rag/ingestion/ & ... & ... \\
%         src/cti\_rag/retrieval/ & Hybrid retrieval & ... \\
%         src/cti\_rag/rag/ & Chain, prompts & ... \\
%         src/cti\_rag/evaluation/ & RAGAS evaluation & ... \\
%         \bottomrule
%     \end{tabular}
%     \caption{Prototype components with path, function, and technical details}
%     \label{tab:components}
% \end{table}

% TODO Abbildung: UI-Screenshot (Chat-Ansicht mit Pipeline-Flow-Card oder
% Landing Page). Druckbarkeit des Dark-Themes pruefen, ggf. helles Theme.
% \begin{figure}[htbp]
%     \centering
%     \includegraphics[width=\textwidth]{ui_chat_view}
%     \caption{Analyst-facing user interface of the prototype}
%     \label{fig:ui-overview}
% \end{figure}

\section{Evaluation design}
\label{sec:evaluation-design}
% EINZIGE Stelle fuer das Evaluationsdesign (frueheres 2.4 ist hier aufgegangen).

\subsection{Test set construction}
\label{sec:test-set}
% TODO: Query-Set (Anzahl, Kategorien, Erstellung, Ground Truth),
% Verweis auf Anhang C.
...

% TODO Tabelle: Query-Set-Uebersicht.
% \begin{table}[htbp]
%     \centering\small
%     \begin{tabular}{lrp{0.5\textwidth}}
%         \toprule
%         \textbf{Task category} & \textbf{Queries} & \textbf{Example query} \\
%         \midrule
%         Vulnerability assessment & ... & ... \\
%         IOC enrichment & ... & ... \\
%         Threat actor context & ... & ... \\
%         \bottomrule
%     \end{tabular}
%     \caption{Evaluation query set by CTI task category}
%     \label{tab:query-set}
% \end{table}

\subsection{RAGAS metrics}
\label{sec:ragas-metrics}
% TODO: verwendete Metriken definieren (faithfulness, answer relevancy,
% context precision/recall, ...), Judge-Modell, Konfiguration.
...

% TODO Tabelle: Metrik-Definitionen kompakt.
% \begin{table}[htbp]
%     \centering\small
%     \begin{tabular}{lp{0.5\textwidth}l}
%         \toprule
%         \textbf{Metric} & \textbf{Measures} & \textbf{Range} \\
%         \midrule
%         Faithfulness & ... & 0--1 \\
%         Answer relevancy & ... & 0--1 \\
%         Context precision & ... & 0--1 \\
%         Context recall & ... & 0--1 \\
%         \bottomrule
%     \end{tabular}
%     \caption{Evaluation metrics used in this thesis}
%     \label{tab:metrics}
% \end{table}

\subsection{Baseline and variant configurations}
\label{sec:baselines}
% TODO: Konfigurationsmatrix; Setup A/B dokumentieren. Kapitel 4.1
% verweist NUR hierher, keine Wiederholung dort.
...

% TODO Tabelle: Konfigurationsmatrix (wichtigste Tabelle von 3.7).
% \begin{table}[htbp]
%     \centering\small
%     \begin{tabular}{llll}
%         \toprule
%         \textbf{ID} & \textbf{Configuration} & \textbf{Retrieval} & \textbf{Generation model} \\
%         \midrule
%         V0 & Non-retrieval baseline & --- & local 8B \\
%         V1 & Sparse-only & BM25 & local 8B \\
%         V2 & Dense-only & vector & local 8B \\
%         V3 & Hybrid & BM25 + vector & local 8B \\
%         V4 & Hybrid + reranking & BM25 + vector + rerank & local 8B \\
%         V5 & Ablation: hosted model & wie V4 & hosted \\
%         \bottomrule
%     \end{tabular}
%     \caption{Baseline and variant configurations}
%     \label{tab:configurations}
% \end{table}

\subsection{Experimental protocol}
\label{sec:protocol}
% TODO: Ablauf, Reproduzierbarkeit (Seeds, Versionen, Hardware),
% Fairness-Bedingungen der Vergleiche.
...
